% Part: lambda-calculus
% Chapter: syntax
% Section: terms

\documentclass[../../../include/open-logic-section]{subfiles}

\begin{document}

\olfileid[id]{lam}{syn}{trm}
\olsection{Suku}

Suku-suku kalkulus lambda dibangun secara induktif dari persediaan variabel
yang tak hingga banyaknya, $\Obj{v_0}$, $\Obj{v_1}$, \dots, simbol
``$\lambd$'', dan tanda kurung. Kita akan menggunakan $x$, $y$, $z$, \dots{}
untuk menyatakan variabel, serta $M$, $N$, $P$, \dots{} untuk menyatakan suku.

\begin{defn}[Suku] \ollabel{defn:term}
Himpunan \emph{suku} kalkulus lambda didefinisikan secara induktif sebagai
berikut:
\begin{enumerate}
  \item \ollabel{defn:term-var} Jika $x$ adalah variabel, maka $x$ adalah
    suku.
  \item \ollabel{defn:term-abs} Jika $x$ adalah variabel dan $M$ adalah
    suku, maka $(\lambd[x][M])$ adalah suku.
  \item \ollabel{defn:term-app} Jika $M$ dan $N$ keduanya suku, maka
    $(MN)$ adalah suku.
\end{enumerate}
\end{defn}

Jika suku $(\lambd[x][M])$ dibentuk menurut
\olref{defn:term-abs}, kita mengatakan bahwa suku itu merupakan hasil suatu
\emph{abstraksi}, dan $x$ dalam $\lambd[x]$ disebut
\emph{!!{parameter}}. Suku $(MN)$ yang dibentuk menurut
\olref{defn:term-app} merupakan hasil suatu \emph{aplikasi}.

Suku-suku yang didefinisikan di atas bertanda kurung lengkap. Hal ini dapat
menjadi cukup merepotkan, sebagaimana diperlihatkan oleh suku
$(\lambd[x][((\lambd[x][x])(\lambd[x][(xx)]))])$. Kita akan memperkenalkan
konvensi untuk menghindari tanda kurung. Namun, definisi resmi tersebut
memudahkan kita menentukan bagaimana suatu suku dibangun menurut
\olref{defn:term}. Sebagai contoh, langkah terakhir dalam membentuk suku
$(\lambd[x][((\lambd[x][x])(\lambd[x][(xx)]))])$ pasti berupa abstraksi
dengan !!{parameter}~$x$. Suku itu diperoleh melalui abstraksi dari suku
$((\lambd[x][x])(\lambd[x][(xx)]))$, yang merupakan aplikasi dua suku.
Masing-masing dari kedua suku ini merupakan hasil abstraksi, dan demikian
seterusnya.

\begin{prob}
Uraikan pembentukan $(\lambd[g][(\lambd[x][(g (x x))])
  (\lambd[x][(g (x x))])])$.
\end{prob}

\end{document}
